% Section 01: Introduction & Operational Context

\section{Introduction \& Operational Context}
\label{sec:introduction}

This report describes the design and implementation of the \textbf{High-Concurrency Vendor Engine}, a software component developed as part of the CS F213 Object Oriented Programming course project at BITS Pilani (Summer Term 2026). The project forms part of Track 4 (Application Development with Java Frameworks), focusing on building robust, high-performance, and resilient digital solutions to address real-world operational bottlenecks.

\subsection{Operational Context}
BITS Pilani's annual cultural festival is a large-scale event serving a user base of 5,000--6,000 active participants. The digital ecosystem supporting the festival facilitates real-time operations, including:
\begin{itemize}
    \item Ticketing and entry validation.
    \item Event scheduling and coordination across 90+ events.
    \item Food ordering at multiple vendor stalls.
    \item Digital wallet transactions totaling over ₹1.5+ crore in value.
    \item Live dashboards displaying aggregate sales, event progress, and crowd distribution.
\end{itemize}

The central backend of this ecosystem is implemented as a Django monolith, supported by PostgreSQL for persistent storage, Redis as an in-memory cache, Celery for asynchronous background task execution, and WebSockets for real-time notifications. The primary user interface for participants is a mobile application.

\subsection{The Bottleneck Challenge}
During peak festival hours, food stalls experience an intense order overload. In this high-throughput scenario, the central Django monolith and its WebSocket handlers frequently encounter performance degradation. The major operational pain points identified include:
\begin{enumerate}
    \item \textbf{WebSocket Connection Drops}: Simultaneous order streams from thousands of users strain the central WebSocket server, leading to connection dropouts and order delivery delays.
    \item \textbf{Stall App Lag}: The vendor-facing application lags due to synchronous network dependency, causing long waiting lines, order processing delays, and incorrect revenue reporting.
    \item \textbf{Lack of Offline Resilience}: In the event of a Wi-Fi or network failure on the festival floor, vendors lose the ability to accept, process, or track orders. This leads to substantial data loss, double-booking, and financial discrepancies.
\end{enumerate}

\subsection{Proposed Solution: The Vendor Engine}
To resolve these bottlenecks, the \textbf{High-Concurrency Vendor Engine} has been developed. It is a lightweight, local desktop application built in Java 17 that runs on-premise at each vendor stall. By moving order routing, prioritization, and queue management to local vendor machines, the system eliminates direct WebSocket dependencies during peak operation. 

Key features of the Vendor Engine include:
\begin{itemize}
    \item \textbf{High-Concurrency Ingestion}: An 8-thread consumer pool capable of parsing and routing up to 50 simulated orders/second.
    \item \textbf{Strict Prioritization}: A custom queue sorting strategy that prioritizes VIP festival-goers and handles tie-breakers chronologically (First-In, First-Out).
    \item \textbf{Autonomous Offline Failover}: An independent network monitoring daemon that detects connectivity drops and serializes the system state locally to prevent data loss.
    \item \textbf{Automatic Synchronization}: On reconnection, the daemon automatically compiles and uploads served orders back to the central monolith, restoring consistency.
\end{itemize}

% Source: archive/Vendor_Engine_Project_Context.md:5-18
% Source: README.md:1-16
